\chapter{KẾT LUẬN VÀ HƯỚNG PHÁT TRIỂN}
\label{ch:conclusion}

\section{Kết luận và Đánh giá mức độ hoàn thành chuẩn CDIO}
\label{sec:cdio_evaluation}
Sau 3 tuần nghiên cứu và triển khai nghiêm túc dưới sự hướng dẫn của TS. Phan Thanh Toàn, nhóm đã hoàn thành toàn diện các mục tiêu đề ra cho Đề tài số 18:
\begin{enumerate}
    \item \textbf{C -- Conceive (Hình thành ý tưởng):} Đã phân tích thấu đáo bài toán ô nhiễm dầu tràn và dầu vón cục tại vùng biển miền Trung Việt Nam; mô hình hóa toán học các hiện tượng vật lý SAR và xác lập tiêu chuẩn phân biệt đốm tối.
    \item \textbf{D -- Design (Thiết kế hệ thống):} Thiết kế thành công kiến trúc đường ống xử lý ảnh 4 bước chuẩn mực, ánh xạ trực tiếp với 4 chương trọng tâm trong giáo trình Xử lý ảnh.
    \item \textbf{I -- Implement (Cài đặt kỹ thuật):} Tự cài đặt 100\% các thuật toán cốt lõi (Lọc trung vị, Phân ngưỡng Otsu tối ưu, Toán tử hình thái học, Trích biên $\beta(A)$, Gradient Sobel) bằng thư viện \texttt{NumPy} từ số 0; kiểm chứng sai số bằng 0 tuyệt đối so với OpenCV.
    \item \textbf{O -- Operate (Vận hành & Đánh giá):} Vận hành thử nghiệm trên toàn bộ 5,630 ảnh SAR CSIRO, đạt Recall 91.23\% và tốc độ xử lý hơn 21 FPS; hoàn thiện ứng dụng Web tương tác Streamlit phục vụ trình diễn trực quan.
\end{enumerate}

\section{Ưu điểm và Hạn chế của Đề tài}
\label{sec:pros_cons}

\subsection{Ưu điểm nổi bật}
\begin{itemize}
    \item \textbf{Minh bạch toán học 100\% (White-box):} Không sử dụng các mô hình học sâu hộp đen, mọi quyết định phát hiện đều có cơ sở vật lý và công thức giải tích rõ ràng.
    \item \textbf{Tốc độ xử lý siêu nhanh:} Tối ưu hóa ma trận NumPy giúp xử lý một ảnh $400 \times 400$ trong chưa đầy 50 ms.
    \item \textbf{Trực quan hóa chuẩn mực:} Đường viền đỏ ôm khít mép ngoài của vệt dầu theo chuẩn lý thuyết hình thái học.
\end{itemize}

\subsection{Hạn chế tồn tại}
\begin{itemize}
    \item Vẫn còn một tỉ lệ báo động giả (False Positive $\approx 13\%$) do các vùng gió lặng cục bộ ven đảo.
    \item Chưa tự động phân tách phần ảnh đất liền nếu vệ tinh chụp sát mép bờ.

\end{itemize}

\section{Định hướng phát triển và Ứng dụng thực tiễn}
\label{sec:future_work}
Nhóm đề xuất 2 hướng nâng cấp mang ý nghĩa thực tiễn to lớn trong giai đoạn tiếp theo:
\begin{enumerate}
    \item \textbf{Mô phỏng quỹ đạo trôi dầu Lagrangian (Lagrangian Oil Drift Simulation):}
    Lấy tọa độ và hình dạng vệt dầu phát hiện được làm điều kiện biên ban đầu, tích hợp trường vector gió bề mặt tái phân tích ERA5 và mô hình hải lưu đại dương (HYCOM/NEMO) theo phương trình vi phân chuyển động hạt:
    \begin{equation}
        \frac{d\vec{x}}{dt} = \vec{u}_{\text{current}} + \alpha \cdot \vec{u}_{\text{wind}}
    \end{equation}
    (với hệ số trôi gió $\alpha \approx 0.03$), giúp dự báo chính xác sau 24--48 giờ vệt dầu sẽ trôi dạt vào bãi biển nào của miền Trung để chủ động phòng chống.
    \item \textbf{Tích hợp dữ liệu định danh vệ tinh AIS truy vết tàu xả thải (AIS Tracking):}
    Liên kết tọa độ và thời gian phát hiện vệt dầu với cơ sở dữ liệu định vị tàu biển quốc tế AIS. Bằng cách đối soát đường đi của các tàu chở dầu đi qua tọa độ đó tại thời điểm vệ tinh Sentinel-1 chụp ảnh, cơ quan chức năng Việt Nam có thể xác định danh tính tàu vi phạm và lập biên bản xử phạt hành chính về tội hủy hoại môi trường biển.
\end{enumerate}
